% !TeX root = ../../main.tex
\section{推广到无穷维}

\begin{solution}[δ函数缩放性质]
\problem{证明 $\delta(ax)=\delta(x)/|a|$。（提示：考虑 $\int \delta(ax)\,d(ax)$，并记住 $\delta(x)=\delta(-x)$。）}

\answer
目标是证明一维δ函数在变量缩放下满足
\[
\delta(ax)=\frac{1}{|a|}\delta(x), \quad a\neq 0.
\]

严格证明采用“对任意测试函数的作用一致性”。设 $g(x)$ 为任意光滑且快速衰减函数，则δ函数的定义性质为
\[
\int_{-\infty}^{\infty}\delta(x)g(x)\,dx=g(0).
\]

现在考虑表达式 $\delta(ax)$ 对 $g(x)$ 的作用：
\[
I=\int_{-\infty}^{\infty}\delta(ax)\,g(x)\,dx.
\]

作变量代换 $y=ax$，则 $x=y/a$，且 $dx=dy/a$。于是
\[
I=\int_{-\infty}^{\infty}\delta(y)\,g(y/a)\,\frac{1}{|a|}\,dy.
\]

这里使用了 $dy/a$ 在积分区间反向时需引入绝对值，故得到 $1/|a|$。

利用$ \delta $函数的抽样性质：
\[
\int \delta(y)\,h(y)\,dy = h(0),
\]
可得
\[
I=\frac{1}{|a|}g(0).
\]

另一方面，如果存在函数 $F(x)$ 满足
\[
F(x)=\frac{1}{|a|}\delta(x),
\]
则其对任意 $g(x)$ 的作用为
\[
\int F(x)g(x)\,dx=\frac{1}{|a|}g(0),
\]
与上式完全一致。

因此在分布意义下必有
\[
\delta(ax)=\frac{1}{|a|}\delta(x).
\]

\end{solution}


\begin{solution}[δ函数复合函数]
\problem{ 证明
\[
\delta(f(x))=\sum_i \frac{\delta(x-x_i)}{\left|\frac{df}{dx}\right|_{x_i}},
\]
其中 $x_i$ 是 $f(x)$ 的零点。提示：考虑 $\delta(f(x))$ 在何处“发散”，并在零点附近对 $f(x)$ 作泰勒展开，保留首个非零项。}

\answer
设 $x_i$ 为 $f(x)$ 的简单零点，即 $f(x_i)=0$ 且 $f'(x_i)\neq 0$。由于$ \delta $函数的支撑只集中在函数为零的位置，因此 $\delta(f(x))$ 只在各零点附近有贡献。

考虑任意测试函数 $g(x)$，研究积分
\[
I=\int \delta(f(x))\,g(x)\,dx.
\]

将积分分解为各零点邻域之和：
\[
I=\sum_i \int_{x\sim x_i} \delta(f(x))\,g(x)\,dx.
\]

在每个零点附近对 $f(x)$ 作泰勒展开，保留首项：
\[
f(x)=f(x_i)+ (x-x_i)f'(x_i)+\cdots \approx (x-x_i)f'(x_i).
\]

因此在该邻域内可作变量替换
\[
y=f(x)\approx f'(x_i)(x-x_i), \quad dy=f'(x_i)\,dx.
\]

于是积分变为
\[
\int_{x\sim x_i} \delta(f(x))g(x)\,dx
=
\int \delta(y)\,g(x(y))\,\frac{dy}{|f'(x_i)|}.
\]

利用$ \delta $函数抽样性质得
\[
=\frac{g(x_i)}{|f'(x_i)|}.
\]

将所有零点贡献相加：
\[
I=\sum_i \frac{g(x_i)}{|f'(x_i)|}.
\]

另一方面，如果存在分布恒等式
\[
\delta(f(x))=\sum_i \frac{\delta(x-x_i)}{|f'(x_i)|},
\]
则对任意 $g(x)$ 作用得到
\[
\int \left(\sum_i \frac{\delta(x-x_i)}{|f'(x_i)|}\right) g(x)\,dx
=
\sum_i \frac{g(x_i)}{|f'(x_i)|},
\]
与上式一致。

因此在分布意义下成立
\[
\delta(f(x))=\sum_i \frac{\delta(x-x_i)}{|f'(x_i)|}.
\]

结论成立，前提是所有零点均为简单零点（即 $f'(x_i)\neq 0$）。
\end{solution}

\begin{solution}[$\theta$ 函数与 $\delta$ 函数关系]
\problem{考虑阶跃函数 $\theta(x-x')$，当 $x-x'$ 为负时函数值为0，当 $x-x'$ 为正时函数值为1。证明
\[
\delta(x-x')=\frac{d}{dx}\theta(x-x').
\]}

\answer
阶跃函数定义为
\[
\theta(x-x')=
\begin{cases}
0, & x<x',\\
1, & x>x'.
\end{cases}
\]

因此它在 $x\neq x'$ 的区域内是常数函数，故有
\[
\frac{d}{dx}\theta(x-x')=0,\quad x\neq x'.
\]

但在 $x=x'$ 处函数发生跳跃（从0跳到1），其导数在经典意义下不存在，但在分布意义下对应$ \delta $函数。

为了严格说明这一点，考虑任意测试函数 $g(x)$，计算
\[
\int_{-\infty}^{\infty}\frac{d}{dx}\theta(x-x')\,g(x)\,dx.
\]

对其进行分部积分：
\[
=\left[\theta(x-x')g(x)\right]_{-\infty}^{\infty}-\int_{-\infty}^{\infty}\theta(x-x')g'(x)\,dx.
\]

由于 $\theta(x-x')=0$ 在 $x\to -\infty$，且 $\theta(x-x')=1$ 在 $x\to +\infty$，边界项为
\[
g(\infty)-0.
\]

更直接地将积分分段：
\[
\int_{x'}^{\infty} g'(x)\,dx = g(\infty)-g(x').
\]

于是得到
\[
\int \frac{d}{dx}\theta(x-x')\,g(x)\,dx = g(x').
\]

另一方面，$ \delta $函数满足定义性质
\[
\int \delta(x-x')g(x)\,dx=g(x').
\]

因此在分布意义下必有
\[
\frac{d}{dx}\theta(x-x')=\delta(x-x').
\]

结论成立，即$ \delta $函数可以视为阶跃函数的导数。
\end{solution}


\begin{solution}[弦的正常模展开（连续系统）]
\problem{在 $t=0$ 时，一根弦的位移为
\[
\psi(x,0)=
\begin{cases}
\frac{2xh}{L}, & 0\le x \le \frac{L}{2},\\[4pt]
\frac{2h}{L}(L-x), & \frac{L}{2}\le x \le L,
\end{cases}
\]
证明其后续时间演化可以展开为
\[
\psi(x,t)=\sum_{m=1}^\infty \sin\!\left(\frac{m\pi x}{L}\right)\cos(\omega_m t)\cdot \frac{8h}{\pi^2 m^2}\sin\!\left(\frac{m\pi}{2}\right).
\]}

\answer
对于两端固定的弦，其本征模为
\[
\phi_m(x)=\sqrt{\frac{2}{L}}\sin\!\left(\frac{m\pi x}{L}\right),
\]
因此任意初态可展开为
\[
\psi(x,t)=\sum_{m=1}^\infty \phi_m(x)\cos(\omega_m t)\langle m|\psi(0)\rangle.
\]

首先计算展开系数
\[
\langle m|\psi(0)\rangle=\sqrt{\frac{2}{L}}\int_0^L \sin\!\left(\frac{m\pi x}{L}\right)\psi(x,0)\,dx.
\]

将积分分段：
\[
=\sqrt{\frac{2}{L}}\left[\int_0^{L/2}\frac{2xh}{L}\sin\!\left(\frac{m\pi x}{L}\right)dx
+\int_{L/2}^L \frac{2h}{L}(L-x)\sin\!\left(\frac{m\pi x}{L}\right)dx\right].
\]

对第一项分部积分（取 $u=x$, $dv=\sin(\frac{m\pi x}{L})dx$），得
\[
\int_0^{L/2} x\sin\!\left(\frac{m\pi x}{L}\right)dx
=
-\frac{Lx}{m\pi}\cos\!\left(\frac{m\pi x}{L}\right)\Big|_0^{L/2}
+\frac{L}{m\pi}\int_0^{L/2}\cos\!\left(\frac{m\pi x}{L}\right)dx.
\]

计算可得
\[
=\frac{L^2}{m^2\pi^2}\sin\!\left(\frac{m\pi}{2}\right)
-\frac{L^2}{2m\pi}\cos\!\left(\frac{m\pi}{2}\right).
\]

第二项同理计算：
\[
\int_{L/2}^L (L-x)\sin\!\left(\frac{m\pi x}{L}\right)dx
=
\frac{L^2}{m^2\pi^2}\sin\!\left(\frac{m\pi}{2}\right)
+\frac{L^2}{2m\pi}\cos\!\left(\frac{m\pi}{2}\right).
\]

两项相加后，余弦项相互抵消，只剩正弦项：
\[
\int_0^L \cdots dx
=
\frac{2L^2}{m^2\pi^2}\sin\!\left(\frac{m\pi}{2}\right).
\]

因此系数为
\[
\langle m|\psi(0)\rangle
=
\sqrt{\frac{2}{L}}\cdot \frac{2h}{L}\cdot \frac{2L^2}{m^2\pi^2}
\sin\!\left(\frac{m\pi}{2}\right)
=
\frac{4h\sqrt{2L}}{m^2\pi^2}\sin\!\left(\frac{m\pi}{2}\right).
\]

代回展开式，并与归一化因子合并，可得最终结果
\[
\psi(x,t)=\sum_{m=1}^\infty \sin\!\left(\frac{m\pi x}{L}\right)\cos(\omega_m t)\cdot
\frac{8h}{\pi^2 m^2}\sin\!\left(\frac{m\pi}{2}\right).
\]

证毕。
\end{solution}
